\section{Spectral-sequence machinery}
\label{chap:spectral-sequence-machinery}

Nested cycle and boundary subobjects determine the internal pages,
differentials, and page passage. The constructions below also associate
internal spectral sequences to bounded filtered complexes. Identifying
two such constructions requires compatible comparison maps; existence of
one construction does not automatically provide the comparisons used in
the Adams and synthetic arguments. The optional Mathlib adapters are
maintained in the legacy SII-MATH/KIP126 repository; this Blueprint records
the internal constructions retained in Kervaire126.

\subsection{Canonical internal pages and morphisms}

\begin{definition}[Internal nested-subobject spectral sequence]
  \label{def:internal-nested-spectral-sequence}
  \lean{KIP126.Core.SpectralSequence.SSData,KIP126.Core.SpectralSequence.PreSS,KIP126.Core.SpectralSequence}
  \uses{def:mathlib-abelian}
  \leanok
  At each degree, an ambient object carries decreasing cycles $Z_r$,
  increasing boundaries $B_r\subseteq Z_r$, and their infinite levels.
  The page is the specified quotient $Z_r/B_r$. A spectral sequence adds
  the starting page and differential degree, actual quotient differentials,
  their square-zero law, and the successor kernel/image identifications.
\end{definition}

\begin{definition}[Internal morphism and its induced page maps]
  \label{def:internal-canonical-page-map}
  \lean{KIP126.Core.SpectralSequence.SSDataMorphism,KIP126.Core.SpectralSequence.SSDataMorphism.cycleMap,KIP126.Core.SpectralSequence.SSDataMorphism.pageMap,KIP126.Core.SpectralSequence.PreSSMorphism,KIP126.Core.SpectralSequence.SpectralSequenceMorphism}
  \uses{def:internal-nested-spectral-sequence}
  \leanok
  A map of ambient objects preserving cycles and boundaries induces its
  unique map on every quotient, including $E_\infty$. An ordinary morphism
  uses equal starting pages and differential degrees and requires these
  induced maps to commute with the differential. Independently chosen page
  maps cannot witness the differential condition.
\end{definition}

\begin{proposition}[Representative criterion for an internal morphism]
  \label{prop:internal-representative-morphism-criterion}
  \lean{KIP126.Core.SpectralSequence.DifferentialMorphismCriterion,KIP126.Core.SpectralSequence.DifferentialMorphismCriterion.toSSDataMorphism,KIP126.Core.SpectralSequence.DifferentialMorphismCriterion.toMorphism,KIP126.Core.SpectralSequence.DifferentialMorphismCriterion.exists_morphism}
  \uses{def:internal-canonical-page-map}
  \leanok
  An ambient map preserving cycle representatives determines boundary
  preservation when representatives of zero map to zero on the target page.
  If the resulting quotient maps satisfy the representative differential
  equation, they assemble into a morphism of the internal spectral sequences.
  The induced page maps are the canonical cokernel maps already fixed by the
  internal morphism interface.
\end{proposition}
\begin{proof}
  Use exactness of the target page cokernel to lift the image of each source
  boundary.  To compare differentials, pull back the two source page
  projections and apply the representative equation to the universal pair of
  lifts; epimorphism cancellation then gives commutation of the canonical
  quotient maps with the page differential.
  \leanok
\end{proof}

\begin{proposition}[Derived internal page and representative calculus]
  \label{prop:internal-page-calculus}
  \lean{KIP126.Challenge2.PageCalculus,KIP126.Challenge2.RepresentativeCalculus,KIP126.Def.Comparison.StageInterfaces.pageCalculus,KIP126.Def.Comparison.StageInterfaces.representativeCalculus}
  \uses{def:internal-nested-spectral-sequence}
  \leanok
  The next page is the homology of the current page differential. Boundaries
  are cycles, and an $E_2$ common representative identifies the same class.
  A specified nonzero differential forces both named $E_2$ classes to have
  nonzero representatives on that page. An arbitrary differential equality
  supplies no nonvanishing conclusion.
\end{proposition}
\begin{proof}
  Use the canonical quotient homology isomorphism, square-zero, and the same
  cycle representatives in the defining internal relations.
  \leanok
\end{proof}

\begin{proposition}[Third-page labels, cycles and incoming images]
  \label{prop:internal-page-three-representatives}
  \lean{KIP126.Core.SpectralSequence.PageThree.exists_representative,KIP126.Core.SpectralSequence.PageThree.reachesPage_iff_d2_eq_zero,KIP126.Core.SpectralSequence.PageThree.isBoundaryBy_iff_mem_d2_range,KIP126.Core.SpectralSequence.PageThree.represents_iff_sub_mem_d2_range,KIP126.Core.SpectralSequence.PageThree.equal_iff_sub_mem_d2_range}
  \uses{def:internal-nested-spectral-sequence,prop:internal-page-calculus}
  \leanok
  For an internal spectral sequence in modules over any ring, every
  $E_3$ element has an $E_2$ label represented by the same cycle.
  If the starting page is at most two, an $E_2$ label reaches $E_3$
  exactly when its actual $d_2$ vanishes. Its cumulative boundary
  condition after $d_2$ is precisely membership in the actual incoming
  $d_2$ image. Given a common representative, another label represents
  that same $E_3$ element exactly when their difference lies in this image.
  The same criterion detects equality of two represented $E_3$ elements.
  All degrees, labels and page elements remain quantified.
\end{proposition}
\begin{proof}
  \leanok
  Use the canonical quotient epimorphisms and the spectral sequence's
  successor kernel and image identities. Transport the subobjects to
  module kernels and ranges; the quotient zero criterion and differences
  of common cycles prove the equivalences. No separate homology
  identification is supplied as a new premise.
\end{proof}

\begin{proposition}[Finite-page representatives and complete bases]
  \label{prop:internal-finite-page-calculus}
  \lean{KIP126.Core.SpectralSequence.FinitePageCalculus.represents_unique,KIP126.Core.SpectralSequence.FinitePageCalculus.represents_d_zero_of_later,KIP126.Core.SpectralSequence.FinitePageCalculus.next_projection_zero_iff_incoming,KIP126.Core.SpectralSequence.FinitePageCalculus.page_generated_from_representatives,KIP126.Core.SpectralSequence.FinitePageCalculus.represents_before,KIP126.Core.SpectralSequence.FinitePageCalculus.next_page_two_of_kernel,KIP126.Core.SpectralSequence.FinitePageCalculus.represents_next_nonzero_of_incoming_zero,KIP126.Core.SpectralSequence.FinitePageCalculus.differential_target_later_zero,KIP126.Core.SpectralSequence.FinitePageCalculus.represents_add_tail,KIP126.Core.SpectralSequence.FinitePageCalculus.sixth_zero_of_second_kernel,KIP126.Core.SpectralSequence.FinitePageCalculus.fifth_frame_of_second_kernel,KIP126.Core.SpectralSequence.FinitePageCalculus.fifth_zero_of_second_kernel,KIP126.Core.SpectralSequence.FinitePageCalculus.fifth_zero_of_all_basis,KIP126.Core.SpectralSequence.FinitePageCalculus.frame_after_second,KIP126.Core.SpectralSequence.FinitePageCalculus.third_zero_of_second_kernel,KIP126.Core.SpectralSequence.FinitePageCalculus.fourth_d_zero_of_second_kernel,KIP126.Core.SpectralSequence.FinitePageCalculus.vx_nonzero_no_incoming,KIP126.Core.SpectralSequence.FinitePageCalculus.pq_rep3zero_boundary,KIP126.Core.SpectralSequence.FinitePageCalculus.represents_sub_tail,KIP126.Core.SpectralSequence.FinitePageCalculus.page_has_representative,KIP126.Core.SpectralSequence.FinitePageCalculus.reaches_from_support,KIP126.Core.SpectralSequence.FinitePageCalculus.d_zero_of_second_kernel_extends,KIP126.Core.SpectralSequence.FinitePageCalculus.reaches_before,KIP126.Core.SpectralSequence.FinitePageCalculus.choice_frame_of_three_basis,KIP126.Core.SpectralSequence.FinitePageCalculus.choice_single_basis_zero_page,KIP126.Core.SpectralSequence.FinitePageCalculus.choice_next_zero_of_single_frame_nonzero_differential,KIP126.Core.SpectralSequence.FinitePageCalculus.d_zero_of_complete_representatives,KIP126.Core.SpectralSequence.FinitePageCalculus.page_eq_sum_of_complete_representatives,KIP126.Core.SpectralSequence.FinitePageCalculus.same_boundaries_projection_eq,KIP126.Core.SpectralSequence.FinitePageCalculus.middle_add,KIP126.Core.SpectralSequence.FinitePageCalculus.middle_if,KIP126.Core.SpectralSequence.FinitePageCalculus.middle_rep2,KIP126.Core.SpectralSequence.FinitePageCalculus.reaches_earlier,KIP126.Core.SpectralSequence.FinitePageCalculus.represents_sub,KIP126.Core.SpectralSequence.FinitePageCalculus.differential_target_cycle,KIP126.Core.SpectralSequence.FinitePageCalculus.represents_two_self,KIP126.Core.SpectralSequence.FinitePageCalculus.represents_later_zero,KIP126.Core.SpectralSequence.FinitePageCalculus.page_has_initial_representative,KIP126.Core.SpectralSequence.FinitePageCalculus.represents_transport_eq,KIP126.Core.SpectralSequence.FinitePageCalculus.represents_bit,KIP126.Core.SpectralSequence.FinitePageCalculus.differential_two_direct,KIP126.Core.SpectralSequence.FinitePageCalculus.f2_expand,KIP126.Core.SpectralSequence.FinitePageCalculus.generated_one,KIP126.Core.SpectralSequence.FinitePageCalculus.generated_two,KIP126.Core.SpectralSequence.FinitePageCalculus.fourth_zero_of_three_basis,KIP126.Core.SpectralSequence.FinitePageCalculus.next_page_zero_of_kernel,KIP126.Core.SpectralSequence.FinitePageCalculus.finite_basis_expansion,KIP126.Core.SpectralSequence.FinitePageCalculus.fourth_frame_of_four_basis,KIP126.Core.SpectralSequence.FinitePageCalculus.represents_next_nonzero_of_incoming_two_cases,KIP126.Core.SpectralSequence.FinitePageCalculus.represents_next_nonzero_of_incoming_zero_at,KIP126.Core.SpectralSequence.FinitePageCalculus.represents_next_nonzero_of_incoming_two_cases_at,KIP126.Core.SpectralSequence.FinitePageCalculus.fourth_frame_of_four_basis_at,KIP126.Core.SpectralSequence.FinitePageCalculus.represents_next_nonzero_of_not_incoming,KIP126.Core.SpectralSequence.FinitePageCalculus.represents_next_nonzero_of_not_incoming_at,KIP126.Core.SpectralSequence.FinitePageCalculus.basis_vector_ne_zero,KIP126.Core.SpectralSequence.FinitePageCalculus.basis_pair_sum_ne_zero}
  \uses{def:internal-nested-spectral-sequence,prop:internal-page-calculus}
  \leanok
  Common initial-page representatives determine a unique class on each later
  page. A complete finite basis generates those classes, including their
  sums and differences. Explicit differential equations and the actual
  incoming boundary subobjects determine kernels, nonvanishing, and page
  vanishing. These are general deductions for the given spectral sequence;
  they assume no selected computation delivery.
\end{proposition}
\begin{proof}
  Use the specified quotient projections, finite basis expansion over
  $\mathbb F_2$, and the canonical successor homology isomorphism. Each
  nonvanishing deduction retains its incoming-boundary hypotheses.
  \leanok
\end{proof}

\begin{proposition}[Naturality of internal pages and representatives]
  \label{prop:internal-page-naturality}
  \lean{KIP126.Core.SpectralSequence.SSDataMorphism.hasDifferential_reindexed,KIP126.Core.SpectralSequence.SSDataMorphism.representsOnPage_reindexed,KIP126.Core.SpectralSequence.SSDataMorphism.pageMap_unique,KIP126.Core.SpectralSequence.SpectralSequenceMorphism.representsOnPage,KIP126.Core.SpectralSequence.SpectralSequenceMorphism.hasDifferential,KIP126.Challenge2.MorphismCalculus,KIP126.Def.Comparison.StageInterfaces.internalNaturality}
  \uses{def:internal-canonical-page-map,prop:internal-page-calculus}
  \leanok
  The induced finite and infinite page maps preserve identity and composition.
  Common representatives also transport through an arbitrary reindexing of
  degrees when the starting pages agree, using the same ambient morphism
  and its canonical quotient maps. No additivity of the reindexing is needed.
  Differential equalities require the actual differential commutation and
  compatibility of the reindexing with the differential-degree translation.
  This applies to affine shifts as well as ordinary degree-preserving maps.
  Ordinary morphisms carry both to those of the target sequence. Preservation of nonzero classes requires an additional
  hypothesis and is not asserted here.
\end{proposition}
\begin{proof}
  Uniqueness of restriction to the cycle subobject and the cokernel universal
  property prove the functorial laws. The same representative squares and the
  actual differential commutation prove the elementwise statements.
  \leanok
\end{proof}

\subsection{Categorical spectral-sequence language}

\begin{definition}[Spectral sequence in an abelian category]
  \label{def:mathlib-spectral-sequence}
  \lean{CategoryTheory.SpectralSequence}
  \mathlibok
  \uses{def:mathlib-abelian,def:mathlib-homological-complex}
  Let $\mathcal C$ be an abelian category, let $\kappa$ be an index type, let
  $c : \mathbb Z \to \operatorname{ComplexShape}(\kappa)$ specify the shape
  of the differential on each page, and fix a starting page $r_0\in\mathbb Z$.
  A spectral sequence consists of a homological complex $E_r$ of shape $c(r)$
  for every $r\ge r_0$, together with, for every $p\in\kappa$, an isomorphism
  \[
    H(E_r)_p \cong (E_{r+1})_p.
  \]
  This is Mathlib's definition \texttt{CategoryTheory.SpectralSequence} in
  the pinned revision.  The project uses this upstream language for its generic tower
  constructions. The stage-boundary statements use
  Definition~\ref{def:internal-nested-spectral-sequence}.
\end{definition}

% SOURCE: Mathlib/Algebra/Homology/SpectralSequence/Basic.lean, declaration CategoryTheory.SpectralSequence.Hom.
\begin{definition}[Mathlib spectral-sequence morphism]
  \label{def:mathlib-spectral-sequence-hom}
  \lean{CategoryTheory.SpectralSequence.Hom}
  \mathlibok
  A Mathlib morphism between spectral sequences has pagewise chain maps
  compatible with passage to homology, but both sequences must have the same
  index type, complex shape, and starting page.  It therefore cannot express
  the classical--synthetic reindexing required later.
\end{definition}

\subsection{Pages, convergence, and comparison}

% SOURCE: Mathlib/CategoryTheory/Triangulated/SpectralObject.lean:20-24 records this bridge as TODO;
% Mathlib/CategoryTheory/Triangulated/HomologicalFunctor.lean:217-240 supplies the three exactness lemmas.


% SOURCE: project bridge required by Mathlib/CategoryTheory/Triangulated/SpectralObject.lean:20-24.


% SOURCE: KIP126/Core/SpectralSequence/HomologicalImage.lean:158-194, homological-image morphism and functor.


% SOURCE: KIP126/Core/Algebra/Completion.lean, quotient-tower and completion-witness interface.
\begin{definition}[Degreewise completion witness]
  \label{def:core-filtration-completion-witness}
  \uses{def:core-filtration,def:core-eventually-zero-filtration,def:core-bounded-above-filtration,def:mathlib-abelian}
  \lean{KIP126.Core.Algebra.Filtration.quotientAt}
  \lean{KIP126.Core.Algebra.Filtration.quotientProjection}
  \lean{KIP126.Core.Algebra.Filtration.quotientTransition}
  \lean{KIP126.Core.Algebra.Filtration.quotientProjection_transition}
  \lean{KIP126.Core.Algebra.Filtration.quotientTransition_comp}
  \lean{KIP126.Core.Algebra.Filtration.quotientTransition_id}
  \lean{KIP126.Core.Algebra.Filtration.quotientTower}
  \lean{KIP126.Core.Algebra.Filtration.quotientTowerCone}
  \lean{KIP126.Core.Algebra.Filtration.CompletionWitness}
  \lean{KIP126.Core.Algebra.Filtration.CompletionWitness.hasLimit}
  \lean{KIP126.Core.Algebra.Filtration.CompletionWitness.completionIso}
  \lean{KIP126.Core.Algebra.Filtration.CompletionWitness.completionIso_hom_comp_limit_π}
  \lean{KIP126.Core.Algebra.Filtration.eq_bot_of_le_of_eq_bot}
  \lean{KIP126.Core.Algebra.Filtration.quotientProjection_isIso_of_eq_bot}
  \lean{KIP126.Core.Algebra.Filtration.quotientTransition_isIso_of_eq_bot}
  \lean{KIP126.Core.Algebra.Filtration.quotientTower_isEventuallyConstantTo_of_eq_bot}
  \lean{KIP126.Core.Algebra.Filtration.quotientTowerCone_isLimit_of_eq_bot}
  \lean{KIP126.Core.Algebra.Filtration.CompletionWitness.of_isEventuallyZero}
  \lean{KIP126.Core.Algebra.Filtration.CompletionWitness.of_isBoundedAbove}
  \leanok
  For a filtered graded object $A$ and a grading degree $i$, the quotient
  tower has value $A_i/F^sA_i$ at $s$ and transition
  $A_i/F^sA_i\to A_i/F^tA_i$ when $t\le s$.  Its canonical cone has point
  $A_i$ and legs the quotient projections.  A degreewise completion witness
  is precisely a proof that this canonical cone is limiting, and it supplies
  existence of the corresponding limit.  With that limit chosen, it gives the
  canonical isomorphism
  \[
    A_i\ \cong\ \varprojlim_s A_i/F^sA_i,
  \]
  characterized by its composites with every quotient projection.  The
  witness is explicit data: this definition does not prove that arbitrary
  filtrations are complete, exhaustive, or convergent.  It does
  prove that a degreewise eventually-zero filtration, and in particular a
  degreewise bounded-above filtration, supplies this completion witness.
\end{definition}

\begin{proposition}[Completion from compatible quotient representatives]
  \label{prop:completion-unique-nat-lifts}
  \uses{def:core-filtration-completion-witness}
  \lean{KIP126.Core.Algebra.Filtration.CompletionWitness.separated,KIP126.Core.Algebra.Filtration.CompletionWitness.of_unique_nat_lifts,KIP126.Core.Algebra.Filtration.CompletionWitness.of_unique_nat_module_lifts}
  \leanok
  For a filtration of abelian groups or modules, suppose every family
  in the nonnegative quotient levels compatible with adjacent transitions
  has a unique lift to the original object. Then the canonical cone to
  all integer quotient levels is limiting. Every completion witness
  implies separatedness: a subobject contained in every filtration level
  is zero.
\end{proposition}
\begin{proof}
  Lift the nonnegative levels elementwise and extend to a negative level
  via its quotient transition from level zero. The forgetful functor
  reflects the resulting limit. For separatedness, all quotient
  projections kill the subobject inclusion; limit extensionality makes
  the inclusion zero.
  \leanok
\end{proof}

\begin{proposition}[Completeness of a residual image filtration]
  \label{prop:residual-image-filtration-complete}
  \lean{KIP126.Core.InverseSequence.exists_limit_of_surjective_sub_shift,KIP126.Core.InverseSequence.eq_zero_of_mem_ranges_of_injective_sub_shift,KIP126.Core.Algebra.Filtration.CompletionWitness.of_residual_sequence}
  \uses{prop:completion-unique-nat-lifts,prop:synthetic-ess-coherent-tower-of-finite}
  \leanok
  Let $p_k:B_k\to A$ commute with the transition maps of an inverse
  sequence, and give $A$ the actual image filtration. If $1-\mathrm{shift}$
  is bijective on $\prod_k B_k$ and each $\ker p_k$ is finite, the
  canonical filtration quotient cone is limiting.
\end{proposition}
\begin{proof}
  Lift successive differences of quotient representatives to $B_k$ and
  solve their difference equation to obtain a common representative in
  $A$. Finite projection fibers provide a compatible lift of an element
  in every filtration image; injectivity of the difference operator
  forces this lift, and hence the element, to vanish. Thus the common
  representative is unique.
  \leanok
\end{proof}

% SOURCE: MainPaper/main.tex:253-271 and 827-862; convergence is a concrete
% witness for the admissible sequences, not a generic Core construction.
\begin{definition}[Complete and separated decreasing filtration]
  \label{def:complete-separated-filtration}
  \uses{def:core-filtration,def:core-filtration-completion-witness}
  \lean{KIP126.Core.SpectralSequence.Filtration.IsExhaustive}
  \lean{KIP126.Core.SpectralSequence.Filtration.IsHausdorff}
  \lean{KIP126.Classical.Adams.IsAdamsFiltrationSeparated}
  \leanok
  Let $F$ be a decreasing filtration by subobjects of a graded object $A$.
  Degreewise ordinary exhaustiveness and separatedness mean, respectively,
  \[
    \bigvee_{s\in\mathbb Z}F^sA_n=\top
    \qquad\text{and}\qquad
    \bigwedge_{s\in\mathbb Z}F^sA_n=\bot
  \]
  for every degree $n$.  For graded abelian groups these are the familiar
  union and intersection conditions.  Degreewise completeness means that the
  canonical map
  \[
    A_n\longrightarrow\varprojlim_s A_n/F^sA_n
  \]
  is an isomorphism for every $n$.  KIP126 records these three obligations
  separately: completeness uses the explicit truncation-tower universal
  property, while exhaustiveness and Hausdorffness are predicates on the
  filtration.  This separation keeps a concrete convergence theorem from
  acquiring any of the three conditions by construction.  These interfaces
  do not identify ordinary
  exhaustiveness with the stronger condition that $F^sA_n=\top$ for all
  sufficiently small $s$, nor separatedness with the stronger condition that
  $F^sA_n=\bot$ for all sufficiently large $s$.  Those stronger eventual
  conditions are recorded separately by the bounded-below and bounded-above
  predicates.  Any later convergence theorem must
  state explicitly whether it assumes the ordinary conditions, these stronger
  eventual conditions, or degreewise boundedness.
\end{definition}

% SOURCE: MainPaper/main.tex:134-150, 264-295, 799-823, and 1277-1317.
\begin{definition}[Spectral-sequence page-level convention]
  \label{def:ss-page-level}
  \lean{KIP126.Core.SpectralSequence.PageLevelConvention}
  \leanok
  A page-level convention consists of a first displayed page, a map from its
  natural-number display index to Mathlib's integer page index, a
  cycle/boundary level, and, when synthetic quotients are involved, a quotient
  exponent.  It must preserve successors and record the lower bounds that make
  every page admissible.  In the AIM Adams convention, for every $r\ge2$,
  these functions have values
  \[
    \operatorname{page}(r)=r,\qquad
    \operatorname{cycleLevel}(r)=r-1,\qquad
    \operatorname{quotientExponent}(r)=r-1.
  \]
  Thus the displayed page is $E_r\cong Z_{r-1}/B_{r-1}$ and the associated
  synthetic quotient is $\nu X/\lambda^{r-1}$.  Equality of the last two
  numerical values is a proved convention law, not a definitional
  identification of cycles with a synthetic quotient.
\end{definition}

% SOURCE: MainPaper/main.tex:134-150 and 799-807, Adams differential degree.
% VERIFIED-IN: KIP126/Classical/Adams/Basic.lean, declarations Bidegree,
% classicalAdamsShift, classicalAdamsTarget, classicalAdamsShape, and
% classicalAdamsShape_two_rel.
\begin{definition}[The classical Adams $E_2\to E_3$ page passage]
  \label{def:classical-adams-e2-e3-shape-slice}
  \uses{def:mathlib-spectral-sequence,def:mathlib-complex-shape-up,def:core-f2-module-category,thm:mathlib-module-category-abelian}
  \lean{KIP126.Classical.Adams.ClassicalAdamsSpectralSequence}
  \lean{KIP126.Classical.Adams.classicalAdamsShape_two_rel}
  \lean{KIP126.Classical.Adams.ClassicalAdamsSS.e₂ToE₃}
  \leanok
  In the cohomological Adams convention, a spectral sequence of
  $\mathbb F_2$-modules indexed by $\mathbb Z\times\mathbb Z$ and displayed
  from $E_2$ has differentials
  \[
    d_r:E_r^{s,t}\longrightarrow E_r^{s+r,t+r-1}.
  \]
  In particular, $d_2$ sends $(s,t)$ to $(s+2,t+1)$, and passage to the
  next page gives the isomorphism
  \[
    E_3^{s,t}\cong
      \ker(d_2:E_2^{s,t}\to E_2^{s+2,t+1})\big/
      \operatorname{im}(d_2:E_2^{s-2,t-1}\to E_2^{s,t}).
  \]
\end{definition}

% SOURCE: MainPaper/main.tex:134-150 and 799-807, Adams differential degree.
\begin{definition}[Classical and synthetic Adams complex shape]
  \label{def:adams-complex-shape}
  \uses{def:mathlib-complex-shape-up,def:ss-page-level}
  \lean{KIP126.Classical.Adams.classicalAdamsShape}
  \lean{KIP126.Classical.Adams.classicalAdamsPageLevel}
  \lean{KIP126.Synthetic.SpectralSequence.syntheticAdamsShape}
  \lean{KIP126.Synthetic.SpectralSequence.syntheticAdamsPageLevel}
  \lean{KIP126.Synthetic.SpectralSequence.syntheticAdamsShape_rel}
  \leanok
  For every integer page $r\ge2$, the classical shape on
  $\mathbb Z\times\mathbb Z$ is $\operatorname{up}'(r,r-1)$, and the
  synthetic shape on $\mathbb Z\times\mathbb Z\times\mathbb Z$ is
  $\operatorname{up}'(r,r-1,0)$.  Consequently a nonzero component can only
  run from $(s,t)$ to $(s+r,t+r-1)$, or from $(s,t,w)$ to
  $(s+r,t+r-1,w)$, respectively.  The classical and synthetic shape functions
  are each packaged with their own page-level convention, including the first
  page, admissible internal page, successor laws, cycle level, and quotient
  exponent.  The synthetic relation theorem verifies the displayed target for
  every natural-number page.
\end{definition}

% SOURCE: MainPaper/main.tex:264-295 and 999-1007, extension differential degree.
\begin{definition}[Extension spectral-sequence complex shape]
  \label{def:ess-complex-shape}
  \uses{def:mathlib-complex-shape-up,def:ss-page-level}
  \lean{KIP126.Classical.ExtensionSS.etaESSShape}
  \lean{KIP126.Classical.ExtensionSS.etaESSShape_rel}
  \lean{KIP126.Classical.ExtensionSS.etaESSPageLevel}
  \lean{KIP126.Classical.ExtensionSS.differentialDegree}
  \leanok
  The first displayed extension page is $n=0$.  On the paper's $(s,t)$
  indexing, its page-$n$ differential has additive degree $(n,n)$ and is
  represented by $\operatorname{up}'(n,n)$.  Any summand tag distinguishing
  the source and target spectra, and any synthetic weight, is an additional
  index whose transition must be stated separately; it is not encoded by the
  pair $(n,n)$.  The classical $\eta$-ESS packages this shape with a page-level
  convention starting at page zero.  The general internal ESS uses indices
  $(s,k)$, starts at page zero, and has differential degree $(n,-1)$: the
  second entry is the two-term-complex column transition, while the fixed
  external stem recovers the paper's second Adams coordinate.  These are two
  descriptions of the same grading convention and must not be conflated with
  the page degrees of either input spectral sequence.
\end{definition}

% SOURCE: MainPaper/main.tex:134-150 and 253-295.
\begin{definition}[Adams bidegree convention]
  \label{def:adams-bidegree}
  \lean{KIP126.Core.SpectralSequence.AdamsPage}
  \lean{KIP126.Core.SpectralSequence.ExtensionLength}
  \lean{KIP126.Core.SpectralSequence.AdamsPage.ofNat}
  \lean{KIP126.Core.SpectralSequence.AdamsPage.ofInt}
  \lean{KIP126.Core.SpectralSequence.AdamsPage.toExtensionLength}
  \lean{KIP126.Core.SpectralSequence.ExtensionLength.toAdamsPage}
  \lean{KIP126.Classical.Adams.Bidegree}
  \lean{KIP126.Classical.Adams.classicalAdamsShift}
  \lean{KIP126.Classical.Adams.classicalAdamsTarget}
  \lean{KIP126.Classical.Adams.classicalAdamsShapeAt}
  \lean{KIP126.Classical.ExtensionSS.differentialDegree}
  \lean{KIP126.Classical.ExtensionSS.etaESSShapeAt}
  \lean{KIP126.Classical.Adams.adamsStem}
  \leanok
  The two coordinates in a classical Adams bidegree $(s,t)$ are Adams
  filtration and internal degree; its topological stem is $t-s$.  A
  page-$r$ classical differential has
  \[
    d_r:E_r^{s,t}\longrightarrow E_r^{s+r,t+r-1},
  \]
  so it raises filtration by $r$ and lowers the stem by one.  A page-$n$
  extension differential has shift $(n,n)$ and therefore preserves the stem.
  The bidegree, stem, and target functions are implemented with two distinct
  nominal parameter types: \texttt{AdamsPage} contains a natural-number page
  together with the proof that it is at least two, whereas
  \texttt{ExtensionLength} contains an arbitrary natural-number length.
  Adams and ESS degree functions accept only their respective parameter
  types.  There is no coercion between them.  The named conversions
  \texttt{AdamsPage.toExtensionLength} and
  \texttt{ExtensionLength.toAdamsPage} are explicit; the latter additionally
  requires a proof that the length is at least two.  Raw integer pages remain
  in categorical chain-complex shapes and are converted explicitly.
\end{definition}

% SOURCE: MainPaper/main.tex:690-739, 799-823, and 1138-1150.
\begin{definition}[Synthetic tridegree and page-level conversion]
  \label{def:synthetic-tridegree}
  \uses{def:adams-bidegree}
  \lean{KIP126.Synthetic.SpectralSequence.Tridegree}
  \lean{KIP126.Synthetic.SpectralSequence.syntheticAdamsShift}
  \lean{KIP126.Synthetic.SpectralSequence.syntheticAdamsTarget}
  \lean{KIP126.Synthetic.SpectralSequence.syntheticAdamsPageLevel}
  \lean{KIP126.Synthetic.SpectralSequence.syntheticAdamsShape}
  \lean{KIP126.Synthetic.SpectralSequence.syntheticAdamsShape_rel}
  \lean{KIP126.Synthetic.SpectralSequence.nuDegree}
  \lean{KIP126.Synthetic.SpectralSequence.forgetWeight}
  \lean{KIP126.Synthetic.SpectralSequence.forgetWeight_nuDegree}
  \lean{KIP126.Synthetic.SpectralSequence.forgetWeight_add_shift}
  \lean{KIP126.Synthetic.SpectralSequence.lambdaDegree}
  \lean{KIP126.Synthetic.SpectralSequence.lambdaTarget}
  \lean{KIP126.Synthetic.SpectralSequence.SyntheticAdamsSS.e₂ToE₃}
  \lean{KIP126.Synthetic.SpectralSequence.SyntheticAdamsSS.d₂}
  \leanok
  A synthetic Adams class is indexed by $(s,t,w)$, where $s$ is Adams
  filtration, $t-s$ is the topological stem, and $w$ is synthetic weight.  A
  synthetic $d_r$ has degree $(r,r-1,0)$.  Multiplication by $\lambda$ has
  degree $(0,0,-1)$.  Under the AIM regrading used in this paper, the degree
  embedding views a classical $E_2$ class of bidegree $(s,t)$ as having
  tridegree $(s,t,t)$.  These grading operations and the internal
  cycle/boundary definition of the synthetic pages are implemented.  The
  $E_2$-homology identification with $E_3$ uses the internal successor-page
  theorem.  They do not by themselves assert that a classical page is a homotopy group of a
  synthetic quotient: the later rigidity input supplies the isomorphism with
  $S/\lambda^{r-1}$, and every use must retain that witness and its source.
\end{definition}

% SOURCE: MainPaper/main.tex:827-862, definitions of Z_r and B_r.
\begin{definition}[Finite-page elementwise presentation]
  \label{def:ss-elementwise-page-presentation}
  \uses{def:mathlib-spectral-sequence,def:adams-bidegree,def:ss-page-level,def:adams-complex-shape}
  \lean{KIP126.Core.SpectralSequence.SSData.page,KIP126.Core.SpectralSequence.SSData.pageπ,KIP126.Classical.Adams.PageRepresentatives.cycles,KIP126.Classical.Adams.PageRepresentatives.boundaries,KIP126.Classical.Adams.PageRepresentatives.isCycle_iff_represents}
  \notready
  An elementwise presentation includes its displayed-page reindexing.  In the
  AIM Adams convention it chooses nested subgroups of the initial page
  $E_2$, with $Z_1=E_2$, $B_1=0$ and
  \[
    B_{r-1}^{s,t}\subseteq B_r^{s,t}\subseteq Z_r^{s,t}
      \subseteq Z_{r-1}^{s,t}\subseteq E_2^{s,t}.
  \]
  For every $r\ge2$ it supplies isomorphisms
  \[
    E_r^{s,t}\cong Z_{r-1}^{s,t}/B_{r-1}^{s,t}.
  \]
  These isomorphisms commute with page passage.  Before quotienting, the
  displayed differential is represented by a homomorphism
  \[
    \widetilde d_r:Z_{r-1}^{s,t}\longrightarrow
    E_2^{s+r,t+r-1}/B_{r-1}^{s+r,t+r-1};
  \]
  it vanishes on $B_{r-1}$ and therefore induces the categorical page
  differential.  Its kernel inside $Z_{r-1}$ is $Z_r$.  The subgroup $B_r$
  is the inverse image in $E_2$ of the image of the incoming differential in
  $E_2/B_{r-1}$, namely the image of
  \[
    \widetilde d_r:Z_{r-1}^{s-r,t-r+1}
      \longrightarrow E_2^{s,t}/B_{r-1}^{s,t}.
  \]
  This proves $B_{r-1}\subseteq B_r\subseteq Z_r$.
  A filtered spectral sequence initialized at $E_0$ must first apply its
  stored page reindexing before using these Adams symbols.  This node records
  only finite pages; it makes no assertion about an abutment or about all
  future pages.

  Lean implements the internal quotient pages and the paper's cycle and
  boundary submodules on the actual Adams $E_2$ page, with the representative
  criterion for cycles.  A coherent comparison from a concrete Mathlib
  spectral sequence to this presentation remains the separate obligation in
  Proposition~\ref{prop:categorical-elementwise-comparison}; the linked
  internal constructions do not establish that comparison.
\end{definition}

% SOURCE: MainPaper/main.tex:134-151, 827-862, and throughout the appendix.
\begin{definition}[Algebraic $E_\infty$ presentation]
  \label{def:ss-einfty-presentation}
  \uses{def:ss-elementwise-page-presentation}
  \lean{KIP126.Core.SpectralSequence.SSData.eInfty}
  \leanok
  Given compatible finite-page data, define
  $Z_\infty=\bigcap_{r\ge1}Z_r$ and
  $B_\infty=\bigcup_{r\ge1}B_r$ in the chosen graded abelian category, and
  define the algebraic limit page by $E_\infty=Z_\infty/B_\infty$.
  This quotient is not an abutment or an associated graded object; every such
  identification is supplied by a separate convergence witness.
\end{definition}

\begin{definition}[Finite survival and permanent classes]
  \label{def:ss-cycle-boundary-survival}
  \uses{def:ss-elementwise-page-presentation}
  \lean{KIP126.Classical.Adams.PageRepresentatives.IsCycle,KIP126.Classical.Adams.PageRepresentatives.IsBoundary,KIP126.Core.SpectralSequence.RepresentsOnPage,KIP126.Core.SpectralSequence.ReachesPage,KIP126.Core.SpectralSequence.SurvivesTo}
  \leanok
  A representative in the initial $E_2$ page is an $r$-cycle when it lies in
  $Z_r$ and an $r$-boundary when it lies in $B_r$.  An $E_r$ class survives
  through page $R$ when it has a representative in $Z_R$ modulo
  $B_{R-1}$.  This finite survival statement makes no claim about pages after
  $R$ and cannot by itself be upgraded to permanence.
  The implementation distinguishes \texttt{ReachesPage}, which permits a
  zero continuation, from \texttt{SurvivesTo}, which requires a nonzero class
  on that page.
\end{definition}

\begin{definition}[Permanent representative and nonzero $E_\infty$ class]
  \label{def:ss-permanent-einfty-class}
  \uses{def:ss-einfty-presentation,def:ss-cycle-boundary-survival}
  \lean{KIP126.Classical.Adams.PageRepresentatives.IsPermanent,KIP126.Core.SpectralSequence.IsPermanentCycle,KIP126.Core.SpectralSequence.IsFinalBoundary,KIP126.Core.SpectralSequence.NonzeroSurvival}
  \leanok
  A representative is permanent when it lies in $Z_\infty$; its limiting
  class is its coset in $E_\infty=Z_\infty/B_\infty$.  That class is nonzero
  exactly when the representative is not in $B_\infty$.
\end{definition}

% SOURCE: MainPaper/main.tex:253-263, PROJECT_BOUNDARY.md external-input boundary.
\begin{definition}[Convergence and detection witness]
  \label{def:ss-convergence-detection-witness}
  \uses{def:core-filtration,def:ss-einfty-presentation}
  \lean{KIP126.Core.SpectralSequence.Convergence}
  \lean{KIP126.Core.SpectralSequence.Detects}
  \lean{KIP126.Core.SpectralSequence.detect_zero}
  \lean{KIP126.Core.SpectralSequence.detect_difference}
  \leanok
  A convergence-and-detection witness for a concrete spectral sequence consists
  of a graded abutment $A_*$, a decreasing filtration $F^sA_n$, a comparison
  $E_\infty\cong\operatorname{gr}^F A$, and a detection relation between
  nonzero $E_\infty$ classes and filtered elements of $A_*$.  It records the
  naturality and representative ambiguity. Existence of such a witness
  for the chosen sequence must be established separately; it does
  not assert a general Boardman theorem, vanishing of a $\lim^1$ obstruction,
  or automatic convergence of arbitrary unbounded filtrations.
\end{definition}

\begin{lemma}[Projection equation for a restricted differential]
  \label{lem:filtered-associated-graded-differential-projection}
  \uses{def:core-associated-graded-differential,def:core-filtered-chain-complex,def:core-to-associated-graded}
  \lean{KIP126.Core.SpectralSequence.FilteredComplex.toAssociatedGraded_comp_associatedGradedDifferential}
  \leanok
  The quotient projection from $F^sC_k$ to
  $\operatorname{gr}^s_F(C_k)$ followed by the associated-graded differential
  equals the chosen restriction of $d_k$ followed by the quotient projection at
  degree $k-1$.
\end{lemma}

\begin{proof}
  \leanok
  Apply the cokernel universal property to the defining factorisation of the
  restricted differential.
\end{proof}

% SOURCE: KIP126/Def/SpectralSequence/Representatives/Proofs.lean:31-64;
% adapted from KIP-infra/SpectralSequence/FilteredComplex.lean at commit 65de864.
\begin{theorem}[Associated-graded differential and filtered representatives]
  \label{thm:filtered-associated-graded-differential-representative}
  \uses{def:core-filtered-chain-complex,def:core-associated-graded-differential,def:core-to-associated-graded,lem:filtered-associated-graded-differential-projection}
  \lean{KIP126.Core.SpectralSequence.FilteredComplex.associatedGradedDifferential_iff_representative}
  \leanok
  Let $T$ be a projective object and let
  $a:T\to\operatorname{gr}^s_F(C_k)$ and
  $b:T\to\operatorname{gr}^s_F(C_{k-1})$ be generalized elements.  Then
  $a$ followed by the associated-graded differential equals $b$ if and only
  if there are a filtration-level representative $x:T\to F^sC_k$ and a
  filtration-level map $dx:T\to F^sC_{k-1}$ such that $x$ projects to $a$,
  $dx$ is the differential of $x$ after inclusion into the underlying
  complex, and $dx$ projects to $b$.
\end{theorem}

\begin{proof}
  \leanok
  Projectivity lifts $a$ through the quotient map.  The filtered differential
  supplies the restricted map $F^sC_k\to F^sC_{k-1}$, and the projection
  equation identifies its quotient with the associated-graded differential.
  Conversely, cancel the monomorphism $F^sC_{k-1}\hookrightarrow C_{k-1}$
  to identify any stated $dx$ with the restricted differential of $x$, then
  use the projection equation.
\end{proof}

% SOURCE: KIP126/Def/SpectralSequence/Representatives/Proofs.lean:86-109;
% adapted from KIP-infra/SpectralSequence/FilteredComplex.lean at commit 65de864.
\begin{theorem}[Cycle criterion by a higher-filtration differential lift]
  \label{thm:filtered-associated-graded-cycle-representative}
  \uses{thm:filtered-associated-graded-differential-representative,def:core-filtered-chain-complex,def:core-to-associated-graded}
  \lean{KIP126.Core.SpectralSequence.FilteredComplex.associatedGradedDifferential_eq_zero_iff_representative}
  \leanok
  For a projective generalized element $a:T\to\operatorname{gr}^s_F(C_k)$,
  its associated-graded differential is zero if and only if $a$ has a
  representative $x:T\to F^sC_k$ whose differential lifts to
  $F^{s+1}C_{k-1}$.  Thus the vanishing condition is expressed by a
  filtration-level cycle lift, not by choosing a representative in the
  quotient without further data.
\end{theorem}

\begin{proof}
  \leanok
  Apply the preceding representative equivalence with the zero generalized
  element.  The zero projection condition is equivalent, by the cokernel
  universal property, to lifting the differential through the next filtration
  inclusion.
\end{proof}

% SOURCE: KIP126/Def/SpectralSequence/Representatives/Proofs.lean:126-164;
% adapted from KIP-infra/SpectralSequence/FilteredComplex.lean at commit 65de864.
\begin{theorem}[Naturality of filtered differential representatives]
  \label{thm:filtered-differential-representative-naturality}
  \uses{def:filtered-chain-map,thm:filtered-chain-map-restricted-differential}
  \lean{KIP126.Core.SpectralSequence.FilteredComplex.Morphism.differentialRepresentative_naturality}
  \leanok
  Let $u:(C,F)\to(D,G)$ be a filtered chain map.  If a map $dx$ at target
  filtration level $t$ represents the differential of $x$ at source level
  $s$, then applying the filtration restrictions of $u$ to $x$ and $dx$
  gives a map at the same respective levels representing the differential of
  the image of $x$.  The source and target levels are arbitrary; choosing
  $t=s+r$ gives the cycle lift and choosing $s=t-r+1$ gives the boundary lift
  used in page comparisons.
\end{theorem}

\begin{proof}
  \leanok
  Postcompose the representative equation with the chain-map component of
  $u$.  The chain-map square and the two filtration-preservation
  factorisations rewrite the result as the target representative equation.
\end{proof}

\begin{definition}[Associated-graded chain complex]
  \label{def:filtered-associated-graded-complex}
  \uses{def:core-associated-graded-object,def:core-associated-graded-differential,thm:core-associated-graded-differential-sq}
  \leanok
  \lean{KIP126.Core.SpectralSequence.FilteredComplex.associatedGradedComplex}
  For a filtered chain complex $(C,F)$ and fixed filtration degree $s$, the
  associated-graded chain complex has degree-$k$ object
  $\operatorname{gr}^s_F(C_k)$ and differential the compiled
  associated-graded differential.  In the displayed consecutive component the
  differential is followed by the canonical index transport induced by
  $(k+1)-1=k$.  The square-zero theorem supplies the chain-complex
  relation; this node packages those existing components as a Mathlib chain
  complex rather than introducing a second differential.
\end{definition}

\begin{theorem}[Displayed differential of the associated-graded complex]
  \label{thm:filtered-associated-graded-complex-d}
  \uses{def:filtered-associated-graded-complex,def:core-associated-graded-differential}
  \lean{KIP126.Core.SpectralSequence.FilteredComplex.associatedGradedComplex_d}
  \leanok
  In consecutive chain degrees, the differential of the associated-graded
  complex is the induced differential followed by the canonical transport
  along $(k+1)-1=k$.
\end{theorem}

\begin{proof}
  \leanok
  Unfold the chain-complex constructor and simplify its allowed differential
  component.
\end{proof}

% SOURCE: KIP126/Core/SpectralSequence/FilteredComplex.lean:641-764, associated-graded maps and functoriality.
\begin{definition}[Associated-graded map of a filtered chain map]
  \label{def:filtered-associated-graded-map}
  \uses{def:filtered-chain-map,def:core-associated-graded-map}
  \lean{KIP126.Core.SpectralSequence.FilteredComplex.Morphism.associatedGradedMap}
  \leanok
  At each filtration and chain degree, a filtered chain map induces the
  associated-graded map obtained from its filtration-level factorisation.
\end{definition}

\begin{theorem}[Naturality of the associated-graded differential]
  \label{thm:filtered-associated-graded-map-comm}
  \uses{def:filtered-chain-map,def:filtered-associated-graded-map,def:core-associated-graded-map,def:core-associated-graded-differential,thm:filtered-chain-map-restricted-differential,lem:core-associated-graded-map-projection,lem:filtered-associated-graded-differential-projection}
  \lean{KIP126.Core.SpectralSequence.FilteredComplex.Morphism.associatedGradedMap_comm}
  \leanok
  For every filtration degree $s$ and chain degree $k$, the map induced by a
  filtered chain map $u$ satisfies
  \[
    \operatorname{gr}^s(u)_k\,;\,d^{\operatorname{gr},G}_{s,k}
      =d^{\operatorname{gr},F}_{s,k}\,;\,\operatorname{gr}^s(u)_{k-1}.
  \]
\end{theorem}

\begin{proof}
  \leanok
  Precompose with the quotient epimorphism from $F^sC_k$.  The universal
  property of the two cokernels replaces the two associated-graded maps by
  their filtration-level factorisations, and the preceding restriction
  theorem identifies the middle composites.  The corresponding projection
  equation for the source differential then finishes by cancellation.
\end{proof}

\begin{definition}[Associated-graded chain map]
  \label{def:filtered-associated-graded-chain-map}
  \uses{def:filtered-chain-map,def:filtered-associated-graded-complex,def:filtered-associated-graded-map,def:core-associated-graded-map}
  \leanok
  \lean{KIP126.Core.SpectralSequence.FilteredComplex.Morphism.associatedGradedChainMap}
  A filtered chain map $u:(C,F)\to(D,G)$ induces a chain map on every fixed
  filtration grade.  Its component in chain degree $k$ is the map
  \[
    \operatorname{gr}^s_F(C_k)\longrightarrow\operatorname{gr}^s_G(D_k),
  \]
  and the chain-map equation is compatibility with the compiled
  associated-graded differentials.
\end{definition}

\begin{definition}[Functorial associated-graded chain maps]
  \label{def:filtered-associated-graded-functoriality}
  \uses{def:filtered-associated-graded-chain-map,def:filtered-associated-graded-map,def:core-filtered-morphism-functoriality}
  \lean{KIP126.Core.SpectralSequence.FilteredComplex.Morphism.associatedGradedMap_id}
  \lean{KIP126.Core.SpectralSequence.FilteredComplex.Morphism.associatedGradedMap_comp}
  \lean{KIP126.Core.SpectralSequence.FilteredComplex.Morphism.associatedGradedChainMap_id}
  \lean{KIP126.Core.SpectralSequence.FilteredComplex.Morphism.associatedGradedChainMap_comp}
  \lean{KIP126.Core.SpectralSequence.FilteredComplex.associatedGradedFunctor}
  \leanok
  The associated-graded construction is functorial in filtered chain maps:
  identities induce identity maps, composites induce composites, and the
  fixed-filtration chain maps assemble into
  \texttt{associatedGradedFunctor}.  These equations are part of the compiled
  interface and are used by the later spectral-object adapter.
\end{definition}

% SOURCE: KIP126/Core/SpectralSequence/Convergence.lean, explicit selected-page/abutment comparison witness.


% SOURCE: MainPaper/main.tex:461-631 and 1218-1235, maps with shifted degrees.
\begin{definition}[Reindexed spectral-sequence morphism]
  \label{def:reindexed-ss-morphism}
  \uses{def:adams-bidegree,def:synthetic-tridegree}
  \lean{KIP126.Comparison.ClassicalSynthetic.ReindexedSpectralSequenceMap}
  \lean{KIP126.Comparison.ClassicalSynthetic.ReindexedSpectralSequenceMap.differential_comm}
  \lean{KIP126.Synthetic.SpectralSequence.fixedWeightPage}
  \lean{KIP126.Comparison.ClassicalSynthetic.forgetWeight_differential_degree}
  \lean{KIP126.Comparison.ClassicalSynthetic.sourcePageMap,KIP126.Comparison.ClassicalSynthetic.sourceEInftyMap}
  \leanok
  For every fixed synthetic weight, the internal classical--synthetic
  comparison supplies a map of ambient objects preserving the nested cycle
  and boundary subobjects.  The maps on finite pages and $E_\infty$ are
  induced from this same map, and the comparison requires those determined
  page maps to commute with the actual differentials.  It does not choose
  independent page-passage isomorphisms.  The degree projection is
  $(s,t,w)\mapsto(s,t)$; binding a comparison to the actual Adams tower and
  the selected synthetic family remains a construction obligation.

  This fixed-weight construction uses the displayed degree projection.
  General affine reindexings require additional grading comparisons.
\end{definition}

% SOURCE: MainPaper/main.tex:138-150, 159-162, 1606-1753.
\begin{definition}[Multiplicative spectral sequence]
  \label{def:multiplicative-ss}
  \uses{def:ss-elementwise-page-presentation}
  \lean{KIP126.Classical.Adams.SphereAdamsAlgebraPresentation,KIP126.Classical.Adams.ExternalAdamsPairingLaws,KIP126.Classical.Adams.SphereAdamsModuleLaws}
  \notready
  A pairing of spectral sequences specifies, on every displayed page,
  products $E_{1,r}^i\otimes E_{2,r}^j\to E_{3,r}^{i+j}$.  A multiplicative
  spectral sequence is, by definition, the self-pairing obtained by taking
  all three spectral sequences to be the same one.  Associativity,
  commutativity, and unitality are additional properties rather than fields
  silently built into this definition.  Passage to the next page is
  multiplicative, and the
  page-$r$ differential is a derivation
  \[
    d_r(xy)=d_r(x)y+(-1)^{|x|}x d_r(y)
  \]
  for an explicitly chosen additive parity function, including for
  multi-index gradings; over $\mathbb F_2$ the sign disappears but the degree
  shift does not.  A pairing of converging spectral sequences additionally
  records the product on $E_\infty$, a filtered abutment product
  $F^sA\,F^{s'}A\subseteq F^{s+s'}A$, the induced associated-graded product,
  and their compatibility under the convergence isomorphism.  Module
  spectral sequences use the analogous action laws and must not be promoted
  to internally multiplicative objects.  This interface is required before
  expressions such as $h_j^2$ or $h_0x$ are meaningful.

  The linked Lean interfaces currently supply the sphere product, external
  pairings and module laws on $E_2$, including the $d_2$ Leibniz rule for the
  product and external pairing.  They do not yet supply this node's products
  on every page, their page-passage compatibility, or the filtered-abutment
  product and convergence comparison.
\end{definition}

% SOURCE: MainPaper/main.tex:253-263 and all later elementwise arguments.


\subsection{Filtered-complex realization}

% SOURCE: KIP126/Core/SpectralSequence/FilteredComplex.lean: filtered-complex morphism declarations; Mathlib/Algebra/Homology/HomotopyCategory/SpectralObject.lean.
\begin{definition}[Filtered chain map]
  \label{def:filtered-chain-map}
  \uses{def:mathlib-chain-complex,def:core-filtered-chain-complex,def:core-filtered-morphism}
  \leanok
  \lean{KIP126.Core.SpectralSequence.FilteredComplex.Morphism}
  \lean{KIP126.Core.SpectralSequence.FilteredComplex.Morphism.toFilteredMorphism}
  \lean{KIP126.Core.SpectralSequence.FilteredComplex.Morphism.id}
  \lean{KIP126.Core.SpectralSequence.FilteredComplex.Morphism.comp}
  For filtered chain complexes $(C,F)$ and $(D,G)$, a filtered chain map is a
  Mathlib chain-complex morphism $u:C\to D$ whose underlying graded morphism is
  filtration preserving.  Equivalently, for every $(s,k)$ it includes the
  restriction $F^sC_k\to G^sD_k$ commuting with the two subobject inclusions;
  the chain-map square makes these restrictions commute with the restricted
  differentials.  Identities and composites retain this data.
\end{definition}

% SOURCE: KIP126/Core/SpectralSequence/FilteredComplex.lean:193-344, fixed-level complexes and filtration inclusion chain maps.
\begin{definition}[Fixed-level filtered complexes and inclusions]
  \label{def:filtered-chain-level-complex-and-inclusions}
  \uses{def:core-filtered-chain-complex,def:core-filtration-inclusions}
  \lean{KIP126.Core.SpectralSequence.FilteredComplex.filteredChainComplex}
  \lean{KIP126.Core.SpectralSequence.FilteredComplex.filteredChainComplex_X}
  \lean{KIP126.Core.SpectralSequence.FilteredComplex.filteredChainComplex_d}
  \lean{KIP126.Core.SpectralSequence.FilteredComplex.filtration_inclusion_comp_differential}
  \lean{KIP126.Core.SpectralSequence.FilteredComplex.filtrationInclusionChainMap}
  \lean{KIP126.Core.SpectralSequence.FilteredComplex.filtrationInclusionChainMap_f}
  \lean{KIP126.Core.SpectralSequence.FilteredComplex.filtrationInclusionChainMap_id}
  \lean{KIP126.Core.SpectralSequence.FilteredComplex.filtrationInclusionChainMap_comp}
  \lean{KIP126.Core.SpectralSequence.FilteredComplex.instCategory}
  \leanok
  At a fixed filtration level, the restricted objects and differentials form a
  genuine chain complex.  Every order relation $t\le s$ gives a chain map
  $F^sC\to F^tC$; its component is the canonical filtration inclusion, and
  these maps satisfy identity and composition laws.  The explicit
  differential-compatibility lemma records why the inclusions are chain maps.
\end{definition}

\begin{definition}[Fixed-level filtered-complex functor]
  \label{def:filtered-chain-level-functor}
  \uses{def:filtered-chain-map,def:core-filtered-chain-complex,def:filtered-chain-level-complex-and-inclusions}
  \lean{KIP126.Core.SpectralSequence.FilteredComplex.Morphism.filteredChainMap}
  \lean{KIP126.Core.SpectralSequence.FilteredComplex.Morphism.filteredChainMap_id}
  \lean{KIP126.Core.SpectralSequence.FilteredComplex.Morphism.filteredChainMap_comp}
  \lean{KIP126.Core.SpectralSequence.FilteredComplex.filteredChainComplexFunctor}
  \lean{KIP126.Core.SpectralSequence.FilteredComplex.Morphism.filteredChainDiagramNatTrans}
  \leanok
  Restricting a filtered chain map to a fixed filtration degree gives a chain
  map between the two restricted complexes.  These restrictions preserve
  identities and composition, and assemble into a natural transformation of
  the decreasing filtration diagrams.
\end{definition}

\begin{theorem}[Compatibility of filtered restrictions]
  \label{thm:filtered-chain-map-restricted-differential}
  \uses{def:filtered-chain-map,def:core-filtered-chain-complex}
  \lean{KIP126.Core.SpectralSequence.FilteredComplex.preserves_differential}
  \leanok
  For a filtered chain map $u:(C,F)\to(D,G)$, write
  $\phi^u_{s,k}:F^sC_k\to G^sD_k$ for its chosen restriction and write
  $\delta^F_{s,k}$ and $\delta^G_{s,k}$ for the restricted differentials.
  Then the restriction squares commute:
  \[
    \phi^u_{s,k}\,;\,\delta^G_{s,k}
      =\delta^F_{s,k}\,;\,\phi^u_{s,k-1}.
  \]
\end{theorem}

\begin{proof}
  \leanok
  Postcompose both sides with the target filtration inclusion.  The two
  factorisation equations and the chain-map square identify the resulting
  composites; cancellation of the monomorphism given by that inclusion gives
  the displayed equality.
\end{proof}

\begin{definition}[Filtered complex as a cochain diagram]
  \label{def:filtered-complex-cochain-diagram}
  \uses{def:mathlib-chain-complex,def:mathlib-cochain-complex,def:core-filtered-chain-complex,def:core-filtration-inclusions,def:filtered-chain-level-complex-and-inclusions,def:filtered-chain-level-functor}
  \lean{KIP126.Core.SpectralSequence.FilteredComplex.filteredChainDiagram}
  \lean{KIP126.Core.SpectralSequence.FilteredComplex.filteredCochainComplex}
  \lean{KIP126.Core.SpectralSequence.FilteredComplex.filteredCochainDiagram}
  \lean{KIP126.Core.SpectralSequence.FilteredComplex.Morphism.filteredCochainDiagramNatTrans}
  \lean{KIP126.Core.SpectralSequence.FilteredComplex.filteredChainDiagramFunctor}
  \lean{KIP126.Core.SpectralSequence.FilteredComplex.filteredCochainDiagramFunctor}
  \leanok
  Let $\mathbb Z_{\mathrm{dec}}$ be the filtration category with a unique arrow
  $s\to t$ exactly when $t\le s$.  A filtered chain complex $(C,F)$ determines
  a functor
  \[
    D_F:\mathbb Z_{\mathrm{dec}}\longrightarrow
      \operatorname{CochainComplex}(\mathcal C,\mathbb Z),
    \qquad D_F(s)^n=F^sC_{-n},
  \]
  where $\mathcal C$ is the ambient abelian category.
  The differential is the restriction of $d_{-n}:C_{-n}\to C_{-n-1}$, and an
  arrow $s\to t$ is sent to the inclusion $F^sC\hookrightarrow F^tC$.  This
  records the reversal $k=-n$ and the decreasing-order convention explicitly;
  the implementation uses Mathlib's canonical chain/cochain equivalence rather
  than silently identifying the two abbreviations.  A filtered chain map
  induces a natural transformation of both the chain and cochain diagrams.
\end{definition}


% SOURCE: KIP126/Core/SpectralSequence/SpectralObjectAdapter.lean:26-342, precomposition morphism and filtered triangulated adapter functor.


% SOURCE: KIP126/Core/SpectralSequence/SpectralObjectAdapter.lean:359-394, homological-image adapter morphisms and functor.


% SOURCE: KIP126/Core/SpectralSequence/Convergence.lean, endpoint-extension and boundary-witness interface.


% SOURCE: KIP126/Core/SpectralSequence/Convergence.lean, endpoint-extended spectral-object and first-page construction.


% SOURCE: KIP126/Core/SpectralSequence/Convergence.lean, endpoint-extension morphisms and induced spectral-object morphisms.


% SOURCE: KIP126/Core/SpectralSequence/SpectralObjectAdapter.lean: project adapter declarations;
% Mathlib/Algebra/Homology/HomotopyCategory/SpectralObject.lean;
% Mathlib/CategoryTheory/Triangulated/SpectralObject.lean;
% Mathlib/Algebra/Homology/HomotopyCategory/HomologicalFunctor.lean;
% Mathlib/Algebra/Homology/HomotopyCategory/ShiftSequence.lean.


% SOURCE: MainPaper/main.tex:253-271, Definition def:ess; standard filtered-complex construction used implicitly.
\begin{theorem}[A bounded filtered complex produces an internal spectral sequence]
  \label{thm:filtered-complex-to-spectral-sequence}
  \lean{KIP126.Core.SpectralSequence.FilteredComplex.toSpectralSequence}
  \uses{def:core-filtered-chain-complex,def:core-bounded-filtration,def:internal-nested-spectral-sequence,thm:filtered-associated-graded-differential-representative}
  \leanok
  Let a filtered chain complex in an abelian category have degreewise bounded
  filtration and a filtration-preserving differential. Its actual nested
  cycle and boundary subobjects, their quotient pages, and the induced
  differentials assemble into an internal SSData spectral sequence, starting
  at page $0$, with differential degree $(r,-1)$. The successor kernel and
  image identifications are proved from those same representatives.
  This bounded construction makes no abutment or convergence assertion;
  applications retain their separately stated convergence hypotheses.
\end{theorem}
\begin{proof}
  Assemble the existing filtered-complex pre-spectral-sequence data.
  The proved differential-square, next-cycle kernel, and next-boundary
  image laws discharge the three laws of the internal construction.
  \leanok
\end{proof}
